\documentclass[10pt,a4paper]{amsart}
\usepackage[margin=19mm]{geometry}
\usepackage{amsmath,amssymb,mathtools}
\usepackage[scr=rsfs,cal=euler]{mathalfa}
\usepackage{xfrac}
\usepackage[unicode,hidelinks,bookmarks=false]{hyperref}
\newcommand{\ie}{i.e.\ }
\newcommand{\cf}{cf.\ }
\newcommand{\resp}{respectively\ }
\newcommand{\Subsection}{\subsection}
\providecommand{\nobreakdash}{\nobreak\hbox{-}}
\setlength{\emergencystretch}{2em}
\multlinegap0pt
\usepackage{bm}
\usepackage{bbm}
\usepackage{dsfont}
\usepackage{xspace}
\usepackage{cancel}
\usepackage{mathtools}
\usepackage{amsthm}
\makeatletter
\@ifundefined{theorem}{\newtheorem{theorem}{Theorem}}{}
\@ifundefined{lemma}{\newtheorem{lemma}[theorem]{Lemma}}{}
\makeatother
\newcommand{\End}{\textnormal{End}}
\newcommand{\bN}{\mathbb{N}}
\newcommand{\bR}{\mathbb{R}}
\newcommand{\id}{\textnormal{Id}}
\newcommand{\pr}{\textnormal{pr}}
\title[On bubbling of Fueter maps]{On bubbling of Fueter maps}
\author{Jacek Rzemieniecki}
\date{Source: arXiv:2608.26803v1 (CC BY 4.0)}
\begin{document}
\maketitle

\noindent\emph{Source and licence.} On bubbling of Fueter maps. Version arXiv:2608.26803v1 (CC BY 4.0), distributed under the Creative Commons Attribution 4.0 International licence, as recorded in the arXiv licence metadata for this exact version and shown on the abstract page https://arxiv.org/abs/2608.26803v1. Full rights notice: licenses/arxiv-2608.26803-CC-BY-4.0.txt.\par\smallskip

\noindent\textit{Editorial scope.} Counted unit: the Lemma whose source label is \texttt{dense image} in the article (source0.tex lines 763--769, source label \texttt{dense image}), delivered together with the article's own complete original proof of it (source0.tex lines 771--817, one \texttt{proof} environment of 47 lines). No mathematics is written, repaired or re-derived: statement and proof are the article's own text line for line. Required context retained uncounted: none. Every definition, display and label of those lines is retained unchanged. Omitted from this package and not redistributed: 1--762, 770--770, 818--978. Nothing inside a retained excerpt is omitted or altered: every retained body is delivered exactly as the article's lines are written, so no omission receipt of any kind accompanies this package. Citations: the retained excerpts carry no citation command at all, so no bibliography entry of the article is transcribed and no reference is redistributed. Reference adaptation: none was needed and none was made; every reference inside the delivered unit resolves inside it, so no unresolved reference or citation is produced. Printed slips kept literal: no printed word, symbol, hypothesis or number of the counted statement or of its proof was altered. Layout and class: the article's own class is \texttt{article} with packages including graphicx, inputenc, fontenc, amsmath, amsthm, amsfonts, amssymb, enumitem, natbib, url, slashed, bbm, tikz-cd, mathtools; the wrapper is the lane's pinned \texttt{amsart} editorial wrapper at 10pt on A4, which restates the theorem-like environments the retained text needs and adds the article's own macro definitions that the retained text needs (\texttt{\textbackslash End}, \texttt{\textbackslash bN}, \texttt{\textbackslash bR}, \texttt{\textbackslash id}, \texttt{\textbackslash pr}), with extra wrapper packages (bm, bbm, dsfont, xspace, cancel, mathtools). The delivered numbering is the unit's own. Assets: no retained span contains a figure environment, a graphics-inclusion command, a PDF-page inclusion, a table or a TikZ picture; no image file is copied into this package, the package declares no asset dependency and no image file is redistributed with it. Licence: version arXiv:2608.26803v1 of this article is distributed under the Creative Commons Attribution 4.0 International licence, as recorded in the arXiv licence metadata for this exact version and shown on the abstract page https://arxiv.org/abs/2608.26803v1. A full rights notice is delivered at licenses/arxiv-2608.26803-CC-BY-4.0.txt. Characters: every retained excerpt is pure ASCII, so the delivered engine, XeLaTeX with its default Latin Modern text font, reports no missing character.
\begin{lemma}\label{dense image}
  For every $f \in W^{k, p}(M, X)$, the operator 
  \begin{equation*}
    B_f: C^m(M \times X, \bR) \to W^{k-1, p}\Gamma(f^* TX): \; B_f(h)=\nabla^{g_X} h(\, \cdot \,, f(\, \cdot \,))
  \end{equation*}
  has dense image.
\end{lemma}

\begin{proof}
  Let $\iota: X \hookrightarrow \bR^N$ be an isometric embedding for some $N$ large enough. This embedding induces a universal projection map
  \begin{equation*}
    \Pi: X \to \End (\bR^N): \; \Pi(y) = \pr_{T_y \iota (T_y X)}
  \end{equation*}
  which is clearly smooth so that 
  \begin{equation*}
    \Pi_f \coloneq \Pi \circ f \in W^{k,p}(M, \End (\bR^N)).
  \end{equation*}
  If $s \in W^{k-1, p}\Gamma(f^*TX)$, then $f^*T\iota(s)\in W^{k-1, p}(M, \bR^N)$. This means that there is a sequence $(a_l)\in C^\infty(M, \bR^N)^\bN$ with 
  \begin{equation}\label{defining property of a_l}
    a_l  \to f^*T\iota(s) \quad \text{ in } W^{k-1, p}(M, \bR^N).
  \end{equation}
  Now, for every $l \in \bN$, define $h_l \in C^\infty(M\times X, \bR)$ by
  \begin{equation*}
    h_l(x, y) \coloneq \langle a_l(x), \iota(y) \rangle_{\bR^N_{\text{std}}}.
  \end{equation*}
  Since the embedding $\iota$ was chosen to be isometric, it holds that for all $x \in M$
  \begin{equation}\label{key property of h_l}
    T_{f(x)} \iota (\nabla^{g_X} h_l (x, f(x))) = \Pi_f (x) \cdot a_l(x).
  \end{equation} 
  Because $k-3/p>2$, the multiplication
  \begin{equation*}
    W^{k, p}(M, \End(\bR^N)) \times W^{k-1, p}(M, \bR^N) \to W^{k-1, p}(M, \bR^N)
  \end{equation*}
  is continuous, so that the combination of \eqref{defining property of a_l} and \eqref{key property of h_l}
  \begin{equation}\label{key convergence 1}
    f^*T \iota \Big( \nabla^{g_X} h_l( \, \cdot \, , f(\, \cdot \, ))\Big) \to f^*T \iota (s) \quad \text{ in } W^{k-1, p}(M, \bR^N).
  \end{equation}
  Since the bundle map 
  \begin{equation*}
    f^*T \iota: f^*TX \to \underline{\bR}^N
  \end{equation*}
  and its metric adjoint
  \begin{equation*}
    (f^*T \iota)^*: \underline{\bR}^N \to f^*TX
  \end{equation*} 
  have $W^{k,p}$ coefficients, by the Sobolev multiplication theorem they induce bounded maps on $W^{k-1,p}$. Now, as
  \begin{equation*}
    (f^*T \iota)^* (f^*T \iota )= \id_{f^*TX},
  \end{equation*} 
  we deduce from \eqref{key convergence 1} that 
  \begin{equation*}
    \nabla^{g_X} h_l( \, \cdot \, , f(\, \cdot \, )) \to s \quad \text{ in } W^{k-1, p}\Gamma(f^*TX),
  \end{equation*}
  proving the claim.
\end{proof}

\end{document}
